\section{The original trivial point is an actual parameter-adic chart test}

\begin{definition}[The actual parameter-adic DVR of an original chart]
\label{def:c0n-chart-parameter-adic-dvr}
\uses{conv:sg-setting,def:c0n-chart-test-points}
\lean{EnvelopeProperty.ProjSubvariety.chartTAdicDVR}
\leanok
For an original chart index $j$ with $x_j\notin I_X$, let
$V_{j,t}$ be the localization of $A_j[t]$ at the original principal prime
ideal $(t)$, embedded in its actual fraction field.
This is a genuine birational DVR essentially of finite type over $k$.
\end{definition}

\begin{lemma}[Its whole original parameter order and torus law]
\label{lem:c0n-chart-parameter-adic-orders}
\uses{def:c0n-chart-parameter-adic-dvr}
\lean{EnvelopeProperty.ProjSubvariety.chartTAdicDVR_ord_X_pow_mul,EnvelopeProperty.ProjSubvariety.chartTAdicDVR_ord,EnvelopeProperty.ProjSubvariety.chartTAdicDVR_ord_C,EnvelopeProperty.ProjSubvariety.chartTAdicDVR_torus}
\leanok
For $Q(0)\ne0$ one has $\operatorname{ord}_{V_{j,t}}(t^mQ)=m$.
Thus the order of any nonzero whole polynomial is its least supported
parameter exponent, every nonzero constant from $A_j$ has order zero,
and
\[
 \operatorname{ord}_{V_{j,t}}\!\left(\sum_\ell f_\ell t^\ell\right)
 =\min_\ell\bigl(\operatorname{ord}_{V_{j,t}}(f_\ell)+\ell\bigr).
\]
The equality includes the zero polynomial with order $+\infty$.
\end{lemma}
\begin{proof}
Since $t$ is a prime element and $t$ does not divide a polynomial with
nonzero constant term, the actual principal-prime localization formula
computes the order of $t^mQ$ as $m$.
Factor the lowest parameter power from a nonzero polynomial; the remaining
factor has nonzero constant term, proving the least-supported-exponent
formula.  A nonzero constant is the case $m=0$.
For a nonzero polynomial all supported coefficients are nonzero constants
in the valuation ring, so every term on the right has order exactly its
parameter exponent.  Their minimum is the least supported exponent,
which proves the torus law.  For zero the support is empty and both sides
are $+\infty$.
\end{proof}

\begin{lemma}[The attached original chart point is the original trivial point]
\label{lem:c0n-chart-parameter-adic-trivial-point}
\uses{lem:c0n-chart-parameter-adic-orders,lem:sg-div-points}
\lean{EnvelopeProperty.ProjSubvariety.chartDivPoint_chartTAdicDVR_eq_triv}
\leanok
The original chart divisorial point $v_{j,V_{j,t},1}$ equals the original
trivial point of $X^{\mathrm{an}}$.
\end{lemma}
\begin{proof}
It suffices to compare both points on homogeneous forms.
A homogeneous form in $I_X$ has value zero at both points.
If a homogeneous form is not in $I_X$, its original dehomogenization in
$A_j$ is nonzero, since $x_j\notin I_X$ and the coordinate ring is a
domain.  That constant has parameter-adic order zero, so the attached
chart point gives the form value one, exactly as the original trivial
point does.  Homogeneous forms determine the original projective
seminorms, proving equality.
\end{proof}

\begin{theorem}[All envelope estimates may be supplied on actual chart DVRs]
\label{thm:c0n-envelope-only-chart-estimates}
\uses{thm:c0n-polynomial-chart-envelope-assembly,lem:c0n-chart-parameter-adic-trivial-point}
\lean{EnvelopeProperty.ProjSubvariety.isPSH_envelope_of_polynomialConeIdealOnlyChartEstimates}
\leanok
In the polynomial chart envelope criterion it suffices to supply the raw
estimate on all genuine original chart data $(j,V,b)$ with $b>0$ and the
actual torus law.  The required estimate at the original trivial point
is then derived internally from these same actual chart data.
\end{theorem}
\begin{proof}
Choose an original chart with $x_j\notin I_X$, which exists for a
projective subvariety.  Apply the supplied estimate to the actual chart
DVR $V_{j,t}$, the positive integer $b=1$, and its proved torus law.
Its attached point is the original trivial point by the preceding lemma,
so this is exactly the required trivial-point estimate.
The polynomial chart envelope theorem now applies.
\end{proof}
